\chapter{Vectores}
Los vectores son herramientas matemáticas que nos ayudan a describir fenómenos físicos cómo la velocidad, aceleración , fuerzas, etc. son muy importantes en física e ingenierías, vamos intentar describirlos en este capítulo y mostrar sus características principales.

\section{cantidades vectoriales}
se dice que una cantidad vectorial es aquella que tiene magnitud y dirección, al contrario de las cantidades escalares que solo cuentan con magnitud pero no con dirección.
\begin{caja_amarilla}[Notación: formas de representar un vector]
	En física se usa la notación $\vec{a}$ para referirse al vector a, también se suele usar mucho la notación de negritas
	$\mathbf{a}$.
	\medspace
	\begin{center}
		\begin{tikzpicture}[>=Stealth, scale=0.8] % Aumenté un poco la escala
			% Rejilla de fondo (primero para que no tape los ejes)
			\draw[thin, gray!15] (-1,-1) grid (5,5);
			
			% Ejes cartesianos más visibles
			\draw[thick, black!75, ->] (-1,0) -- (5,0) node[right] {$x$};
			\draw[thick, black!75, ->] (0,-1) -- (0,5) node[above] {$y$};
			
			% Proyecciones (líneas punteadas)
			\draw[dashed, gray!60] (3,0) -- (3,2);
			\draw[dashed, gray!60] (0,2) -- (3,2);
			
			% El vector (Azul Marino y más grueso)
			\draw[ultra thick, navyblue, ->] (0,0) -- (3,2) node[midway, above left] {$\mathbf{v}$};
			
			% Punto A (Origen)
			\filldraw[black] (0,0) circle (2pt) node[anchor=north east] {$(0,0)$};
			
			% Punto B (Extremo/Punta)
			\filldraw[black] (3,2) circle (2pt) node[anchor=south west] {$(3,2)$};
			
		\end{tikzpicture}
	\end{center}
	Un vector se puede representar por medio de coordenadas, por ejemplo:
	 $$ \mathbf{a} = (3,2) $$
	también se puede representar por medio de vectores unitarios:
	$$ \mathbf{a} = 3\hat{x} + 2\hat{y}  $$
	las coordenadas individuales se denominan \textbf{componentes} del vector.
	\medspace
	En general, se representa un vector $\mathbf{v}$ como:
	$$ \mathbf{v}=(v_1,v_2) $$
\end{caja_amarilla}
\subsection*{Suma de Vectores}
La suma de vectores se define sumando componente por componente, es decir, sean los vectores $\mathbf{v}=(v_1,v_2)$,$\mathbf{u}=(u_1,u_2)$, la suma es:

\begin{equation}
	\mathbf{v} + \mathbf{u} = (v_1 + u_1, v_2 + u_2)
\end{equation}
\begin{figure}[h] % [h] significa 'here' (aquí mismo)
	\centering
	\begin{tikzpicture}[>=Stealth, scale=1.0]
		% ... (aquí va todo tu código del dibujo) ...
		\draw[thin, gray!15] (-1,-1) grid (6,5);
		\draw[thick, black!75, ->] (-1,0) -- (6,0) node[right] {$x$};
		\draw[thick, black!75, ->] (0,-1) -- (0,5) node[above] {$y$};
		\coordinate (O) at (0,0);
		\coordinate (V) at (3,1);
		\coordinate (U) at (2,3);
		\coordinate (S) at (5,4);
		\draw[dashed, gray] (V) -- (S);
		\draw[dashed, gray] (U) -- (S);
		\draw[ultra thick, navyblue, ->] (O) -- (V) node[midway, below right] {$\mathbf{v}$};
		\draw[ultra thick, blue!60!cyan, ->] (O) -- (U) node[midway, left] {$\mathbf{u}$};
		\draw[ultra thick, darkred, ->] (O) -- (S) node[midway, above left] {$\mathbf{v} + \mathbf{u}$};
		\filldraw[black] (V) circle (1.5pt);
		\filldraw[black] (U) circle (1.5pt);
		\filldraw[darkred] (S) circle (2pt) node[anchor=south west] {$(v_1+u_1, v_2+u_2)$};
	\end{tikzpicture}
	\caption{Visualización gráfica de la suma de vectores mediante el método del paralelogramo.}
	\label{fig:suma_vectores}
\end{figure}
\subsection*{Multiplicación por un escalar}
Sea $\mathbf{v}=(v_1,v_2)$ un vector y $\alpha$ un escalar, entonces, se define la multiplicación por un escalar de la siguiente manera:
\begin{equation}
	c\mathbf{v}=c(v_1,v_2)=(cv_1,cv_2)
\end{equation}
\begin{figure}[htbp]
	\centering
	\begin{tikzpicture}[>=Stealth, scale=1.0]
		% Rejilla
		\draw[thin, gray!15] (-3,-3) grid (5,3);
		
		% Ejes
		\draw[thick, black!75, ->] (-3,0) -- (5,0) node[right] {$x$};
		\draw[thick, black!75, ->] (0,-3) -- (0,3) node[above] {$y$};
		
		% Coordenadas
		\coordinate (O) at (0,0);
		\coordinate (V) at (2,1);    % Vector original v
		\coordinate (V2) at (4,2);   % 2v (Escalado positivo)
		
		% Vector original v
		\draw[ultra thick, navyblue, ->] (O) -- (V) node[midway, above left] {$\mathbf{v}$};
		
		% Vector 2v (estirado)
		\draw[thick, darkred, ->] (O) -- (V2) node[anchor=south west] {$c\mathbf{v}$};
		
		% Puntos para marcar los extremos
		\filldraw[navyblue] (V) circle (1.5pt);
		\filldraw[darkred] (V2) circle (1.5pt);
	\end{tikzpicture}
	\caption{Representación de la multiplicación de un vector $\mathbf{v}$ por un escalar $c > 1$.}
	\label{fig:multiplicacion_escalar}
\end{figure}

\subsection{Vectores en $\mathbb{R}^n$ }
Vamos a generalizar un poco, el conjunto de todos los pares ordenados se le conoce como $\mathbb{R}^2$, al conjunto de todos los vectores con 3 componentes es denomina $\mathbb{R}^3$. En general $\mathbb{R}^n$ se define como el conjunto de n-tuplas de números reales ordenados $(x_1,x_2,...,x_n)$ donde $x_i\in\mathbb{R}$. Las definiciones de suma vectorial y multiplicación por un escalarse extienden hasta $\mathbb{R}^n$.
\subsection{Propiedades Algebraicas de los Vectores}
\begin{caja_amarilla}[Propiedades de los Vectores en $\mathbb{R}^n$ ]
	Sean $\mathbf{u}, \mathbf{v}$ y $\mathbf{w}$ vectores en $\mathbb{R}^n$ y sean $\alpha$ y $\beta$ escalares. Entonces:
	
	\begin{enumerate}
		\item $\mathbf{u} + \mathbf{v} = \mathbf{v} + \mathbf{u}$
		\item $(\mathbf{u} + \mathbf{v}) + \mathbf{w} = \mathbf{u} + (\mathbf{v} + \mathbf{w})$
		\item $\mathbf{u} + \mathbf{0} = \mathbf{u}$
		\item $\mathbf{u} + (-\mathbf{u}) = \mathbf{0}$
		\item $\alpha(\mathbf{u} + \mathbf{v}) = \alpha\mathbf{u} + \alpha\mathbf{v}$
		\item $(\alpha + \beta)\mathbf{u} = \alpha\mathbf{u} + \beta\mathbf{u}$
		\item $\alpha(\beta\mathbf{u}) = (\alpha\beta)\mathbf{u}$
		\item $1\mathbf{u} = \mathbf{u}$
	\end{enumerate}
\end{caja_amarilla}

